% Auto-generated IEEE LaTeX Benchmark Tables for AVSA
\begin{table}[htbp]
\caption{Computational Latency and Resource Allocation Across Multi-Agent Architecture (NVIDIA RTX 4050 6GB)}
\label{tab:latency_benchmarks}
\centering
\begin{tabular}{lcccc}
\hline
\textbf{Agent / Component} & \textbf{Model / Engine} & \textbf{Precision} & \textbf{Latency (ms)} & \textbf{Throughput (FPS)} \\
\hline
Vision Agent & YOLOv8m & FP16 & 214.96 & 4.7 \\
Motion Agent & SEA-RAFT-M & FP32 & 312.83 & 3.2 \\
Behavior + Reasoning & Dynamic Rule Engine & CPU & $\approx 1.20$ & $>800$ \\
\hline
\textbf{End-to-End Pipeline} & \textbf{Multi-Agent} & \textbf{Hybrid} & \textbf{297.3} & \textbf{3.4} \\
\textbf{Effective with Skip=3} & \textbf{Temporal Subsample} & \textbf{--} & \textbf{--} & \textbf{10.1} \\
\hline
\multicolumn{5}{l}{\footnotesize Peak VRAM Utilization: 4.28 GB} \\
\hline
\end{tabular}
\end{table}

\begin{table}[htbp]
\caption{Ablation Study: Impact of Sustained Counter Threshold ($K$) on Anomaly False Positive Suppression}
\label{tab:sustained_counter_ablation}
\centering
\begin{tabular}{ccccc}
\hline
\textbf{Sustain Window ($K$)} & \textbf{Total Triggers} & \textbf{True Positives} & \textbf{False Positives} & \textbf{Precision} \\
\hline
1 & 120 & 90 & 30 & 0.7500 \\
3 & 29 & 29 & 0 & 1.0000 \\
5 & 18 & 18 & 0 & 1.0000 \\
8 & 11 & 11 & 0 & 1.0000 \\
\hline
\multicolumn{5}{p{0.45\textwidth}}{\footnotesize Note: $K=1$ corresponds to traditional unbuffered thresholding, while $K=5$ corresponds to the proposed AVSA hysteresis debouncer.} \\
\hline
\end{tabular}
\end{table}
